\chapter*{Tóm tắt}
\label{summary}

Kiểm thử tự động website là một hướng tiếp cận quan trọng nhằm giảm công sức kiểm thử lặp lại và phát hiện lỗi hồi quy sớm hơn trong quá trình phát triển phần mềm. Tuy nhiên, trong nhiều quy trình hiện nay, tester vẫn phải tốn nhiều thời gian để phân tích yêu cầu, xây dựng test case, chuẩn bị test data, viết automation script và bảo trì kịch bản khi giao diện thay đổi. Bên cạnh đó, nếu toàn bộ quá trình kiểm thử hồi quy đều phụ thuộc vào AI agent hoặc mô hình ngôn ngữ lớn để suy luận lại từ đầu thì chi phí sử dụng mô hình và thời gian thực thi có thể tăng lên.

Khóa luận này xây dựng một hệ thống hỗ trợ kiểm thử tự động website theo hướng kết hợp mô hình ngôn ngữ lớn, AI agent, browser automation và replay script. Hệ thống cho phép người dùng tạo project kiểm thử, nhập yêu cầu bằng ngôn ngữ tự nhiên, sinh test case và test data, xác nhận test case, chạy test trên trình duyệt, ghi nhận evidence, sinh replay script và xem report sau khi thực thi. Trong kiểm thử hồi quy, người dùng có thể chủ động chạy lại các kịch bản đã ổn định bằng replay script thay vì để agent suy luận lại toàn bộ luồng thao tác.

Hệ thống được đánh giá thông qua case study trên website OrangeHRM và khảo sát trải nghiệm người dùng. Kết quả thực nghiệm cho thấy 64 batch sinh test case đều tạo được ít nhất một test case có cấu trúc hợp lệ, với thời gian sinh trung bình 3.66 giây mỗi batch. Trên 167 lần thực thi, hệ thống đạt tỷ lệ Pass tổng thể 83.2\%; trong đó Agent initial run đạt 110/121 lần Pass, tương ứng 90.91\%, còn replay script đạt 63.04\%. Các phiên Agent initial run thành công là cơ sở để sinh replay script cho những kịch bản có thể tái sử dụng. Khi xét riêng các lần replay thành công, thời gian thực thi trung bình thấp hơn agent thành công khoảng 14.65\%.

Về mức sử dụng LLM, 83 lần chạy agent thuộc phạm vi phân tích token sử dụng tổng cộng 5,349,788 tokens, trong khi 46 lần chạy replay script có dữ liệu thực thi và agent token bằng 0. Nếu các lần replay này được chạy lại bằng agent theo mức token trung bình đã ghi nhận, lượng token có thể tránh sử dụng được ước tính khoảng 2.96 triệu tokens. Kết quả này cho thấy hệ thống phù hợp với mục tiêu giảm công sức tạo test case, test data và script ban đầu, đồng thời tạo cơ sở để tái sử dụng kịch bản kiểm thử và giảm mức phụ thuộc vào LLM trong các lần kiểm thử hồi quy đối với các luồng nghiệp vụ đã ổn định.
